\section{考察と利用範囲}
\subsection{コードの有用性が現れる場面}
第一に、発光の機構と観測の機構を同じ計算で追跡できる。ある線の不足を見たとき、励起率、準安定密度、消光、分岐、光路、透過率のどこを変更したかが明確になる。生成・損失収支と成分別出力は、経験的な係数調整だけで説明を終えず、仮定の影響を調べる手段になる。

第二に、物理項をデータで改訂できる。断面積の新版、未観測の放射分岐、追加の消光相手を、species packや反応集合の変更として扱える。原子データを直す作業と、励起の閉包を変更する作業を分けられるため、複数ガスで同じAr参照を再利用しやすい。関数形が変わればplugin、線形の前提が変わればsolver変更、と作業の範囲も判断しやすい。

第三に、観測設計を定量的に検討できる。分岐線を増やすこと、別上準位を追加すること、別種を追加すること、視線を追加することは、それぞれ増える情報が違う。Ar/Cl$_2$の例では、四本のAr線でも上準位が二つなら温度・密度分離が弱い場合がある。測定Jacobianとprofileを使うことで、誤差を減らすべき装置条件と、独立に測るべきプラズマ量を検討できる。

これらは、適用範囲を明示した分光診断の試作、既存CRMの縮約版の比較、校正と視線条件の感度調査、共同研究でのデータ・反応集合の交換に適している。放電の電力・流量から最終的な組成や加工速度を直接予測する課題では、外部の化学・輸送・表面モデルが別に必要になる。

\subsection{外部実測による検証へ進む条件}
現在の外部候補には、Ar-CCPの線比資料\cite{arellano}と、低圧N$_2$/O$_2$のNO紫外発光資料\cite{schuecke}がある。前者には電子分布、準安定・cascade、断面積の整合性、後者には主要なN$_2$(A)密度など、必要入力の閉包が残る。参考論文中のCRMやPIC/MCCが実験と一致したことは、別の反応集合を持つOESCRの検証ではない。

実測へ進める単位は、gas名一般ではなく、反応集合・遷移・圧力範囲・観測幾何を指定した一つのケースである。校正済みスペクトル、波長応答、LSF、光路、ガス温度、供給粒子密度とその不確かさをまとめ、温度または電子密度の独立参照を対応づける。fitに使わない波長・条件で予測を確認し、合否値は結果を見る前に定める。

BOLSIG+等の固定出力とのEEDF・速度係数比較は、分布の読込と積分の外部数値確認に有用である\cite{hagelaar,lxcat}。ただし、現在のtemplateと実行準備の存在を、独立実行結果との一致と取り違えない。その一致が得られても、光学校正や準位閉包を含む実験的Te精度を単独では保証しない。

Fullerらのactinometry\cite{fuller}やNF$_3$洗浄の分光研究\cite{an2023}は、対象種と実験条件の設計に有益である。一方、同じOESから導出された電子温度やactinometry密度を、そのOESモデルの独立真値として再利用すると循環評価になる。図のdigitization誤差、実験誤差、モデル不一致は区別して扱う必要がある。

\subsection{主要な近似と改善の優先順位}
物理的な限界は、定常・選択準位・線形衝突収支、外部密度の固定、領域間輸送の省略、有効分子帯、入力光学厚さによる簡略捕獲、有効壁損失に集約できる。逆問題の限界は、低次元EEDF族への条件づけ、局所SVDと局所profile、曲率に基づく不確かさである。これらは異なる原因なので、単に自由パラメータを増やして同時に解消できるわけではない。

改善では、まず一つの校正済み実測ケースについて供給密度と励起経路を閉じ、独立診断と保留スペクトルで誤差を測る。その結果に基づいて、準安定励起、cascade、消光、捕獲の必要性を判断する。一般のEEDF再構成や複雑な化学連成を先に増やすと、観測が支えない自由度を増やす可能性がある。現行の監査機能は、こうした拡張の前に何が不足しているかを見定めるためにも有用である。

\subsection{結論}
OESCRは、明示した電子状態と粒子密度から、選択準位の線形CR、発光、視線幾何、装置応答、測定残差までを一貫して計算する。物理項の所在と単位が分かれ、複数種のデータを合成できる点は、モデルの点検と改訂を容易にする。再現例は、この計算経路と条件付き逆推定の機能を支持する。同時にAr/Cl$_2$の回復失敗と多視線例の残る誤差は、スペクトル適合だけで物理量を確定できないことを示した。第三者が利用する際は、指定した反応・入力・観測範囲に対する診断モデルとして扱い、実験的な定量主張を独立検証のある範囲へ限定することが妥当である。
